\section{The reduction instance}\label{sec:instance}

\begin{theorem}[Graph-to-packing reduction]\label{thm:reduction}
Fix an integer $c\ge0$ and a real constant $\rho>0$. There is a
deterministic polynomial-time construction with the following property.
Its input is an explicitly listed simple graph $G=(V,E)$, with
$V=\{1,\ldots,n\}$ and $m=|E|\ge1$, and an integer $0\le k\le n$.
Its output is an integer $B$ and an explicit list of $5B$ rational item
sizes in $(1/6,1)$ such that:
\begin{enumerate}
 \item if $G$ has a vertex cover of size at most $k$, the items pack into
       $B$ unit-capacity bins;
 \item any packing into at most $B+c$ bins yields a vertex cover of size
       at most $k+\rho n$.
\end{enumerate}
The running-time and output-length bounds are in ordinary binary
encoding; their constants may depend on $c$ and $\rho$.
\end{theorem}

We give the construction and its elementary properties in this section.
Propositions~\ref{prop:completeness} and~\ref{prop:soundness} prove the
two implications of Theorem~\ref{thm:reduction}, respectively. Vertex
listing is part of the input representation, so the input length is at
least $n$.

\subsection{Parameters and interval geometry}

The integer $P_0$ below bounds a fixed score used to enforce the bin
patterns. The allowance $K$ will bound items lost from those patterns
when $c$ extra bins are available. Choose, once for the fixed pair
$(c,\rho)$,
\begin{equation}\label{eq:parameters}
 \begin{gathered}
  P_0=100^{11},\qquad K=25(c+1)(1+5P_0),\\
  d\in\N,\qquad d\ge\max\{72K,120K/\rho\}.
 \end{gathered}
\end{equation}
Apply Lemma~\ref{lem:competing-trees} to this $d$, obtaining
$T^+,T^-$, common leaf depth $L\ge1$, and combined depth quotas
$p_1,\ldots,p_L$. Put $P=\sum_{\ell=1}^L p_\ell$ and choose an integer
$D_*\ge1$ bounding the number of children of any node in either tree.
These quantities depend only on $c,\rho$. Set
\[
 R=100(P+nK+1).
\]
The $R$ repetitions of each incidence will make the later edge demand
larger than both the local losses $P$ and the global losses proportional
to $nK$.

Order the graph edges as $e=1,\ldots,m$. A \emph{position} is a pair
$r=(e,j)$ with $1\le j\le R$; positions are ordered lexicographically.
Define
\begin{equation}\label{eq:position-coordinates}
 \begin{aligned}
  \theta_e&=(R+1)^{e-1}-1,
  &H_{e,j}&=j(R+1)^{e-1}-1,\\
  g&=\frac{1}{16(R+1)^m},
  &b_{(e,j)}&=\frac12+gH_{e,j},\qquad \beta_e=g\theta_e.
 \end{aligned}
\end{equation}
Write $H_r=H_{e,j}$ when $r=(e,j)$. For each vertex $v$, let
\[
 \mathcal R_v=\{(e,j):v\text{ is an endpoint of }e,\ 1\le j\le R\}.
\]

Each vertex label $v$ uses its own copy of the following geometry.
The roots of $T^+$ and $T^-$ are represented by the closed side cells
$[0,1]$ and $[4,5]$. A node at depth $\ell\ge1$ has a closed cell of
width
\[
 \lambda_\ell=g[100(D_*+1)]^{-\ell}.
\]
Within a parent with left endpoint $a$, the cell of its $j$th child at
depth $\ell$ is
$[a+4j\lambda_\ell,a+(4j+1)\lambda_\ell]$.
Finally put $\Delta=\lambda_L/10$. For each $r\in\mathcal R_v$, add the
closed \emph{job cell} $[b_r-\Delta,b_r]$ on the plus side.
The \emph{test intervals} are the leaf cells of both trees and all the
job cells, each taken left-closed and right-open. Denote their union by
$\mathcal I_v$. Embedding cells remain closed whenever their geometric
intersection or containment is discussed.

\begin{lemma}[Geometry and coordinate order]\label{lem:geometry}
The embedding has the following properties.
\begin{enumerate}
 \item Every child cell lies inside its parent. Every nonroot cell lies
       in the first tenth of its root cell. Distinct depth-$\ell$ cells
       across the two trees have gaps at least $3\lambda_\ell$.
 \item The numbers $b_r$ belong to $[1/2,9/16)$, and distinct job cells
       have gaps at least $g-\Delta$. All the test intervals of one
       label are pairwise disjoint.
 \item For $1\le\ell\le L$,
       $\Delta\le\lambda_\ell/10$ and
       $\lambda_1+2\Delta<g$. An interval of length at most
       $\lambda_\ell+\Delta$ meets at most one depth-$\ell$ cell. If it
       contains a point of a leaf cell, that depth-$\ell$ cell is the
       leaf's ancestor. An interval of length at most
       $\lambda_1+\Delta$ meets at most one job cell.
 \item For $r=(e,j)$, the two independent order bounds are
       \begin{equation}\label{eq:coordinate-order}
        H_{r'}-H_r\ge\theta_e+1\quad(r'>r),\qquad
        \theta_f-\theta_e\ge H_{e,R}+1\quad(f>e).
       \end{equation}
\end{enumerate}
\end{lemma}

\begin{proof}
The right endpoint of the last child is at most
$(4D_*+1)\lambda_\ell$ beyond its parent's left endpoint. For
$\ell>1$ this is less than $\lambda_{\ell-1}$, since
$4D_*+1<100(D_*+1)$. For $\ell=1$ it is less than $g/25<1/10$.
Sibling gaps are $3\lambda_\ell$. Cells with different parents retain
at least the gap between those parents, which is at least
$3\lambda_{\ell-1}>3\lambda_\ell$; the two root cells are separated
by $3$. Induction proves the first assertion, including all cousin
cells, and descendants remain in their depth-one cells.

Successive $H$ values within edge $e$ differ by $(R+1)^{e-1}$. The
difference from $(e,R)$ to $(e+1,1)$ is the same number. Thus the first
step after a position of edge $e$ is $\theta_e+1$, proving the first
bound in Equation~\eqref{eq:coordinate-order}. Similarly,
\[
 \theta_{e+1}-\theta_e=R(R+1)^{e-1}=H_{e,R}+1,
\]
and later edges only increase this difference. Also
$0\le H_r\le R(R+1)^{m-1}-1<(R+1)^m$, proving the range of $b_r$.
Consecutive baselines differ by at least $g$, so job-cell gaps are at
least $g-\Delta$.

The widths decrease with depth, giving
$\Delta\le\lambda_\ell/10$. Since $100(D_*+1)\ge200$,
\[
 \lambda_1+2\Delta\le\frac65\lambda_1
 \le\frac{3g}{500}<g.
\]
In particular $\Delta<g\le1/16$. Job cells therefore lie strictly to
the right of $1/10$ and strictly to the left of $1$; they meet neither
tree's nonroot cells. Leaf cells have common depth $L$ and positive
gaps, proving disjointness of all tests. Finally,
$\lambda_\ell+\Delta\le(11/10)\lambda_\ell<3\lambda_\ell$, and
$\lambda_1+\Delta<g-\Delta$. These strict inequalities and the stated
gaps prove the intersection bounds. A point of a leaf belongs to its
ancestor cell at every depth, so the uniquely intersected cell at depth
$\ell$ must be that ancestor.
\end{proof}

The geometric separation has two uses. A local interval of length
$\lambda_\ell$, even with the additional shortening $\Delta$, can
intersect only one tree cell at depth $\ell$. The two coordinate-order
bounds will separately control which row positions and which edge
resources can be matched. We next realize these constraints through
five-item bins.

\subsection{Patterns and the complete item inventory}

Items have five roles: rows $X$, anchors $A$, global auxiliaries $U$,
local auxiliaries $D$, and flags $F$. Their subclasses are listed in
Table~\ref{tab:patterns}. A \emph{tuple} consists of five distinct
physical item copies. A \emph{table tuple} has one of the displayed
subclass patterns. Its kind is $\mathrm M$ for the first three rows of
Table~\ref{tab:patterns} and $\mathrm G$ for the last row. Only the roles $X,A,D$
carry vertex labels. A table tuple is \emph{good} when these three
labels are equal.

\begin{table}[htbp]
 \centering
 \begin{tabular}{c c c c c c}
  \toprule
  Kind & $X$ & $A$ & $U$ & $D$ & $F$\\
  \midrule
  $\mathrm M$ & $\Tp$ & $\Zanc$ & $\Up$ & $\DM$ & $\FT$\\
  $\mathrm M$ & $\Tm$ & $\Zanc$ & $\Um$ & $\DM$ & $\FT$\\
  $\mathrm M$ & $\Srow$ & $\Zanc$ & $\Up$ & $\DM$ & $\FM$\\
  $\mathrm G$ & $\Srow$ & $\Yanc$ & $\Waux$ & $\DG$ & $\FG$\\
  \bottomrule
 \end{tabular}
 \caption{The four table patterns, with one subclass in each role.
 Equal vertex labels on $X,A,D$ are a separate requirement for a good tuple.}
 \label{tab:patterns}
\end{table}

The $\mathrm M$ patterns will compare a row's completion with an
anchor deadline; the $\mathrm G$ pattern will match a job row to an
anchor key. The inventory below gives every role the same total count
and gives the three labeled roles the same count at each vertex.
We specify every item copy and a rational coordinate $w(i)$; the actual
size will be defined afterwards. The symbols $b,z$ and the auxiliary
lengths below are attributes of individual items.

\paragraph{Rows.}
For each label $v$ and each node of either tree, \emph{including its
root}, create $d$ rows with baseline $b$ equal to the cell's right
endpoint and coordinate $w=b$. Their subclass is $\Tp$ on the plus
side and $\Tm$ on the minus side. Add $P$ further $\Tp$ rows with
baseline $1$. Thus the tree-row counts are
\begin{equation}\label{eq:tree-row-counts}
 t_{v,+}=d|V(T^+)|+P,\qquad
 t_{v,-}=d|V(T^-)|,\qquad
 t_v=t_{v,+}+t_{v,-}\ge P+2d.
\end{equation}
For each $r\in\mathcal R_v$, create $2d$ rows of subclass $\Srow$,
all with baseline $b_r$. Of these, $d$ are \emph{short}, with
$w=b_r-\Delta$, and $d$ are \emph{long}, with $w=b_r$.
All these rows have label $v$. An intended packing sends $d$ rows at
each position to kind $\mathrm M$ and the other $d$ to kind $\mathrm G$.
Put
\[
 J_v=d|\mathcal R_v|=dR\deg(v),
\]
the number of job rows assigned to each kind in such a packing.

\paragraph{Anchors.}
For label $v$, form a baseline multiset $\mathcal B_v$ consisting of
every tree-row baseline, including padding, and $d$ copies of $b_r$
for each $r\in\mathcal R_v$. The latter are hypothetical baselines,
so $|\mathcal B_v|=t_v+J_v$, rather than the total row count.
For each test interval, replace $d$ copies of its right endpoint by
its left endpoint. Leaf-row copies and the hypothetical job copies
provide these replacements; the different tests use distinct designated
copies. For each resulting deadline $z$, create a $\Zanc$ anchor of
label $v$ with $w=-z$.

Set $\ind\{E\}=1$ when condition $E$ holds, and $0$ otherwise.
For $a\in\R$, let $B_v(a)$ be the number of members of $\mathcal B_v$
at most $a$. Moving an endpoint from $r$ to $l<r$ changes its prefix
indicator by
$\ind\{l\le a\}-\ind\{r\le a\}=\ind\{l\le a<r\}$.
Since the tests are disjoint, their deadline count satisfies exactly
\begin{equation}\label{eq:deadline-prefix}
 \#\{\Zanc\text{ anchors of label }v\text{ with }z\le a\}
 =B_v(a)+d\,\ind\{a\in\mathcal I_v\}.
\end{equation}
There are also $d$ \emph{keys} of subclass $\Yanc$ for every position
$r=(e,j)\in\mathcal R_v$, with label $v$ and coordinate
$w=-b_r-\beta_e$. Their total number is $J_v$.

\paragraph{Global auxiliaries.}
These have no vertex label. Create
$\sum_v(t_{v,+}+J_v)$ items of subclass $\Up$, exactly $dk$ of length
$1$ and the rest of length $0$. Create $\sum_v t_{v,-}$ items of
subclass $\Um$, exactly $d(n-k)$ of length $1$ and the rest of length
$0$. An item of either subclass has coordinate $w=-\len(i)$.
For each edge $e$, create $dR$ \emph{permits} of subclass $\Waux$ with
$w=\beta_e$, and $dR$ \emph{nonpermits} of that same subclass with
$w=\beta_e+\Delta$.

\paragraph{Local auxiliaries.}
For each label $v$, create $t_v+J_v$ items of subclass $\DM$.
For each $\ell=1,\ldots,L$, exactly $p_\ell$ of them have length
$\lambda_\ell$; the remaining $t_v+J_v-P$ have length $0$.
Their coordinate is $w=-\len(i)$. Also create $J_v$ items of subclass
$\DG$, with label $v$ and coordinate $0$.

\paragraph{Flags.}
Create $\sum_vt_v$ items of subclass $\FT$ and $\sum_vJ_v$ each of
subclasses $\FM$ and $\FG$. They have no vertex label and coordinate
$0$.

All multiplicities are nonnegative. In particular,
Equation~\eqref{eq:tree-row-counts} gives
$t_v+J_v-P\ge2d+J_v$ zero-length local items, while the $d$ root rows
on each side guarantee enough global items for the specified length-one
stocks, including when $k=0$ or $k=n$. Define
\begin{equation}\label{eq:inventory-total}
 B=\sum_{v=1}^n(t_v+2J_v).
\end{equation}
Each role has exactly $B$ items. This is immediate for $X,A,D$ from
their labelwise counts. For the other roles, use
$\sum_vJ_v=2dmR$: the $\Waux$ inventory is exactly $\sum_vJ_v$, and
the stated global and flag counts each sum to $B$ as well. The three
labeled roles have equal inventories $t_v+2J_v$ at every label $v$.
All coordinates satisfy $|w(i)|\le6$: cell endpoints lie in $[0,5]$,
lengths are at most $1$, and $0\le\beta_e<1/16$.

\subsection{Encoding the patterns and labels in rational sizes}

Large-base aggregation of bounded integer equations is a classical
arithmetic device; see Anthonisse~\cite[Theorem~6 and Section~5]{Anthonisse1973}.
A recent use in bin-packing constructions appears in
Jansen, Pirotton, and Tutas~\cite[Section~3]{JansenPirottonTutas2025}.
We give the bounded five-item version explicitly and separately control
the defects allowed by extra bins.

The encoding has three tasks: keep every item above $1/6$, control
departures from Table~\ref{tab:patterns}, and turn good-tuple feasibility
into an inequality on the coordinates $w(i)$. Feasibility alone will
allow some departures from the table. Conservation of the total scores
will bound their number when at most $c$ extra bins are used.

For a collection of items, let $N_s$ denote the count of subclass $s$
and $N_r$ the count of role $r$. Consider the following ten homogeneous
count equations: the five role equations
$5N_r-\sum_sN_s=0$, in the order $r=X,A,U,D,F$, followed by
\begin{equation}\label{eq:pattern-count-equations}
 \begin{aligned}
  N_{\Tp}+N_{\Tm}-N_{\FT}&=0,\\
  N_{\Zanc}-N_{\FT}-N_{\FM}&=0,\\
  N_{\Up}+N_{\Um}-N_{\FT}-N_{\FM}&=0,\\
  N_{\DM}-N_{\FT}-N_{\FM}&=0,\\
  N_{\Tm}-N_{\Um}&=0.
 \end{aligned}
\end{equation}
Here $\sum_s$ is over all thirteen subclasses. Let $\delta_j(i)$,
$0\le j<10$, be item $i$'s coefficient in these successive left sides,
and define its integer primary score by
\[
 p(i)=\sum_{j=0}^9 100^j\delta_j(i).
\]
The coefficients depend only on the subclass. Each has absolute value
at most $4$, so
$|p(i)|\le4\sum_{j=0}^9 100^j<P_0$.

For vertex labels set $Q_v=3^v$. The secondary score is $q(i)=Q_v$ for
an $X$ or $D$ item of label $v$, $q(i)=-2Q_v$ for an $A$ item of that
label, and $q(i)=0$ on the other roles. Give item $i$ the actual size
\begin{equation}\label{eq:item-size}
 \begin{gathered}
  a_i=\frac15+\gamma p(i)+\mu q(i)+\nu w(i),\\
  \gamma=\frac1{1000P_0},\qquad
  \mu=\frac{\gamma}{1000Q_n},\qquad
  \nu=\frac{\mu}{1000}.
 \end{gathered}
\end{equation}
Roles and labels specify the construction and remain available as
metadata for its analysis. The output consists of the rational sizes of
the indexed physical copies; the argument requires no recovery of this
metadata from a size and no numerical distinctness of different kinds
of items.

\begin{lemma}[Scalar encoding]\label{lem:encoding}
Every item size belongs to $(1/6,1)$, so every feasible bin has at most
five items. For a five-item tuple, zero total primary score is equivalent
to following Table~\ref{tab:patterns}. A feasible five-item tuple has
nonpositive total primary score; if that score is zero, its total
secondary score is nonpositive. If both totals are zero, feasibility is
equivalent to $\sum_iw(i)\le0$. The complete item inventory satisfies
$\sum_i p(i)=\sum_iq(i)=0$.
\end{lemma}

\begin{proof}
On a five-item tuple, every role-equation total belongs to $[-5,20]$
and each remaining total belongs to $[-5,5]$. All ten are integers of
absolute value less than $100$. If their radix-$100$ combination is
zero, reducing modulo $100$ makes the first total zero: it is a
multiple of $100$ strictly between $-100$ and $100$. Divide by $100$
and repeat to obtain all ten equations. This argument applies to signed
totals as well. The converse is immediate.

The role equations now give exactly one item of each role. With flag
$\FT$, the first subclass equation forces a tree row, the next three
force $\Zanc$, one of $\Up,\Um$, and $\DM$, and the last equation
pairs $\Tm$ precisely with $\Um$. With flag $\FM$, the row is
$\Srow$, the anchor and local item are $\Zanc,\DM$, and the global
item must be $\Up$. With flag $\FG$, the remaining subclasses are
forced to be $\Srow,\Yanc,\Waux,\DG$. These are exactly the four
patterns, each of which satisfies every equation.

The perturbation of one item from $1/5$ is bounded by
\[
 |a_i-1/5|
 \le\frac1{1000}+\frac{2}{10^6P_0}
       +\frac{6}{10^9P_0Q_n}<\frac1{30}.
\]
Thus $a_i>1/6$ and $a_i<7/30<1$. For five items, the total absolute
contribution below the primary scale is at most
$10Q_n\mu+30\nu$, and
\[
 \frac{10Q_n\mu+30\nu}{\gamma}
 =\frac1{100}+\frac{3}{100000Q_n}<1,
 \qquad \frac{30\nu}{\mu}=\frac3{100}<1.
\]
An integer primary total at least $1$ would therefore make the bin's
size exceed $1$. If the primary total is zero, the same argument at
the secondary scale excludes a positive secondary total. When both
vanish, the five copies of $1/5$ already sum to $1$, leaving exactly
the asserted coordinate condition.

Finally, the whole inventory satisfies all ten count equations. Every
role has $B$ items; the tree-row and $\FT$ counts both equal
$\sum_vt_v$; each of the $\Zanc$, combined $\Up,\Um$, and $\DM$
counts equals $\sum_v(t_v+J_v)$; and the $\Tm$ and $\Um$ counts
agree. Thus the primary total is zero. At each label the $X,A,D$
inventories are equal, so its secondary contribution is
$(1-2+1)Q_v(t_v+2J_v)=0$.
\end{proof}

\paragraph{Polynomial construction.}
For fixed $c,\rho$, the trees and all their parameters are finite
constants. In particular $t_v=t$ is a fixed constant, while
$R=O(n+1)$ and
$B=nt+4dmR$. Hence all positions and all item copies can be enumerated
in polynomial time in the explicit graph input. The denominators
\begin{equation}\label{eq:encoding-denominators}
 \begin{aligned}
  D_{\rm geom}&=160(R+1)^m[100(D_*+1)]^L,\\
  D_{\rm size}&=10^9P_0 3^n D_{\rm geom}
 \end{aligned}
\end{equation}
are common denominators for all coordinates $w(i)$ and all sizes $a_i$,
respectively. Indeed, $D_{\rm geom}$ accommodates $g$, every
$\lambda_\ell$, $\Delta$, all nested endpoints, and every baseline and
$\beta_e$; Equation~\eqref{eq:item-size} then gives the second claim.
Their bit lengths are
$O(n+m\log(n+2))$, with constants depending on the fixed parameters.
Since $0<a_i<1$, numerators over $D_{\rm size}$ have the same bound.
Integer powering, arithmetic, item enumeration, and optional fraction
reduction therefore take polynomial bit time. The choice of $d$ and the
finite tree construction are fixed for each reduction; $\rho$ is not a
real-valued runtime input.

\begin{lemma}[Good anchors under an additive allowance]\label{lem:good-anchors}
In any packing of the constructed instance into at most $B+c$ bins,
fewer than $K$ items lie outside five-item table tuples. At each vertex
label $v$, at most $3K$ anchors lie outside good tuples. This is a joint
bound on its $\Zanc$ and $\Yanc$ anchors, and therefore applies to any
subset of that label's anchors.
\end{lemma}

\begin{proof}
If the packing has $N\le B+c$ bins, five-item capacity gives $N\ge B$.
The number of unused five-item slots is $5(N-B)\le5c$. Hence there
are at most $5c$ nonfull bins and at most $25c$ items in them. All full
bins have nonpositive integer primary totals by
Lemma~\ref{lem:encoding}. The whole primary total is zero, so the
sum of the magnitudes of the negative full-bin totals is at most
$25cP_0$, supplied by the nonfull items. There are consequently at most
$25cP_0$ negative-primary full bins. The number of items outside
primary-zero full bins is at most
\begin{equation}\label{eq:primary-loss}
 25c+5(25cP_0)=25c(1+5P_0)<K.
\end{equation}
Those primary-zero bins are precisely the table tuples.

In a table tuple write $v_X,v_D,v_A$ for its three labels. Its
secondary score is
$Q_{v_X}+Q_{v_D}-2Q_{v_A}\le0$. If $v_X>v_A$, then
$Q_{v_X}\ge3Q_{v_A}$, already exceeding the negative contribution;
the same argument applies to $v_D$. Thus
\[
 v_X\le v_A,\qquad v_D\le v_A.
\]
Equality of the secondary total occurs exactly when all three labels
are equal.

Here is the exact conservation law that controls smaller-label imports.
Fix a prefix $S=\{1,\ldots,h\}$ and a role $r\in\{X,D\}$.
Let $O_A(S)$ and $O_r(S)$ count items of the indicated role and labels
in $S$ outside table tuples. Let $E_r(S)$ count table tuples whose
role-$r$ item has label in $S$ and whose anchor has label outside $S$.
The original inventories of either role in $S$ equal
$\sum_{v\in S}(t_v+2J_v)$. Every table anchor in $S$ consumes a
role-$r$ item in $S$. Subtracting this consumption from the available
role-$r$ items gives
\begin{equation}\label{eq:label-prefix-conservation}
 E_r(S)=O_A(S)-O_r(S)\le O_A(S)\le K.
\end{equation}
This identity also holds for the empty prefix.

For anchors of label $v$, any unequal row label must be below $v$ and
is counted by $E_X(\{1,\ldots,v-1\})$. At most $K$ table anchors
of label $v$ have this defect. At most another $K$ have a smaller
local-auxiliary label. At most $K$ anchors of the label lie outside
table tuples altogether. Their union has at most $3K$ members, proving
the joint assertion. Summing this bound over labels permits a global
loss of $3nK$.
\end{proof}

\subsection{Feasibility of good tuples}

Lemma~\ref{lem:good-anchors} now lets us analyze an arbitrary packing
through its good tuples, losing at most $3K$ anchors at each vertex.
The loss is independent of the number of job positions. For good
tuples, the remaining coordinate inequality has the following concrete
meaning.

\begin{lemma}[Completion and key inequalities]\label{lem:good-feasibility}
A good tuple has zero primary and secondary totals. Its feasibility is
therefore described exactly as follows.
\begin{enumerate}
 \item For kind $\mathrm M$, with row baseline $b$, deadline $z$,
       and associated global and local items $U,D$, define
       \begin{equation}\label{eq:completion-coordinate}
        h=b-\Delta\ind\{\text{row is short}\}-\len(U)-\len(D).
       \end{equation}
       The tuple is feasible if and only if $h\le z$, and $h\le b$.
       For any collection of such row--auxiliary assignments and any
       $a\in\R$,
       \begin{equation}\label{eq:completion-prefix}
        \#\{h\le a\}=\#\{b\le a\}+\#\{h\le a<b\}.
       \end{equation}
       The last term counts coverage by the half-open intervals $[h,b)$.
 \item For kind $\mathrm G$, let its key have position $r=(e,j)$,
       its row have position $r'$, and its global item belong to edge
       $f$. Feasibility is equivalent to
       \begin{equation}\label{eq:key-feasibility}
        b_{r'}-b_r+\beta_f-\beta_e
        +\Delta\bigl(\ind\{\text{nonpermit}\}
                      -\ind\{\text{row is short}\}\bigr)\le0.
       \end{equation}
       It implies both $r'\le r$ and $f\le e$. For an exact-position
       long row, the global item must be a permit of edge $e$ or an
       item of an earlier edge. At the exact position and edge, a
       short row admits either a permit or a nonpermit, whereas a long
       row admits only a permit.
\end{enumerate}
\end{lemma}

\begin{proof}
The table pattern gives zero primary score; equal labels give
$Q_v+Q_v-2Q_v=0$. Apply Lemma~\ref{lem:encoding}. In kind
$\mathrm M$, the five coordinates sum to $h-z$, proving the first
criterion. All lengths and the shortening term are nonnegative, so
$h\le b$. The events $\{b\le a\}$ and $\{h\le a<b\}$ partition
$\{h\le a\}$, proving Equation~\eqref{eq:completion-prefix}, including
the endpoints and the empty intervals when $h=b$.

In kind $\mathrm G$, summing the row, key, global, local, and flag
coordinates gives exactly the left side of
Equation~\eqref{eq:key-feasibility}. If $r'>r$, then
Equation~\eqref{eq:coordinate-order} gives a position increase at least
$g(\theta_e+1)$. Since $\theta_f\ge0$, the $\beta$ difference is at
least $-g\theta_e$. The final perturbation is at least $-\Delta$.
The total is therefore at least $g-\Delta>0$, excluding $r'>r$
without any condition on $f$.

Independently, if $f>e$, the $\beta$ increase is at least
$g(H_{e,R}+1)$. Since $H_{r'}\ge0$ and $H_r\le H_{e,R}$, the
position difference is at least $-gH_{e,R}$. Again the total is at
least $g-\Delta>0$, excluding $f>e$ without any condition on $r'$.
For an exact-position long row, a same-edge nonpermit would leave
the strictly positive sum $\Delta$, proving the resource restriction.
For an exact-position, same-edge tuple, the sum is simply
$\Delta(\ind\{\text{nonpermit}\}-
\ind\{\text{row is short}\})$, which gives the remaining assertions.
\end{proof}

For later geometric use, an $\mathrm M$ interval with global length
zero and local length $\lambda_\ell$ has length at most
$\lambda_\ell+\Delta$. With both lengths zero it is empty, except
for a short $\Srow$ row, whose interval is its job test interval.
Thus the intersection bounds of Lemma~\ref{lem:geometry} apply even
when the row's baseline is unrelated to the depth of its local item.
Finally, a plus row has $b\le1$, while a minus row has $b\ge4$ and
is never short. Subtracting a global length at most $1$ and a local
length at most $\lambda_1<1/10$ leaves its completion greater than
$1$. Consequently intervals on either side can cover test points
only on that side.
